\chapter{Quasiexemples i contraexemples}\label{cap:contraexemples}

Una definició s'entén de debò quan se sap què \emph{no} la compleix. Aquest
apèndix recull, capítol per capítol, els exemples que mostren on falla una
afirmació plausible, quina hipòtesi és imprescindible o per què una
construcció no es pot simplificar. Molts provenen de la teoria clàssica de
conjunts o d'altres estructures matemàtiques: no són models dels axiomes del
llibre, sinó contrastos que fan visible el paper de cada axioma.

\section*{Aplicacions, parts i valors de veritat (capítols 1--3)}

\begin{description}
\item[Un conjunt sense punts.] El conjunt inicial \(\mathbf0\) és un conjunt
de ple dret i no té cap element global
(capítol~\ref{cap:aplicacions}). \enquote{Ser un conjunt} i \enquote{tenir
elements} són coses diferents.

\item[La separació global no és autodual.] El dual de l'axioma de separació
global (axioma~\ref{ax:separacio-globals}) diria que dues aplicacions
paral·leles coincideixen si coincideixen després de compondre-les amb
totes les aplicacions cap a \(\mathbf0\); com que d'un conjunt no inicial no
en surt cap, l'enunciat faria iguals totes les aplicacions paral·leles. A
partir d'aquest axioma la dualitat deixa de valer
(observació~\ref{obs:dualitat}).

\item[Un bimorfisme que no és isomorfisme.] En els anells, la inclusió
\(\Z\hookrightarrow\mathbb Q\) és mono i epi però no té inversa. En aquesta
teoria no pot passar: el teorema del bimorfisme
(proposició~\ref{prop:bimorfisme-isomorfisme}) usa l'axioma de
representació, que els anells no tenen.

\item[Dues parts amb el mateix domini.] Els dos punts
\(0,1\colon\mathbf1\hookrightarrow\mathbf2\) són parts de \(\mathbf2\) amb el
mateix domini i no són equivalents: una part és una fletxa, no el seu
domini.

\item[Una injecció de suma sense la retracció del cograf.] L'argument del
cograf (proposició~\ref{prop:cograf-retraccio}) fa mono la injecció
\(i_0\colon A\to A+B\) quan hi ha alguna \(B\to A\). Per a \(A=\mathbf0\) i
\(B=\mathbf1\) no n'hi ha cap; \(i_0\) és mono per una altra raó (tota
aplicació amb domini \(\mathbf0\) ho és), i en general cal la disjunció de
les sumes.

\item[Un element generalitzat repartit.] Siguin \(A=X=\mathbf2\), \(m\) el
punt \(1\colon\mathbf1\hookrightarrow\mathbf2\) i \(x=\id_{\mathbf2}\). Aleshores
\(x\) no pertany a \(m\) ni al seu complement: sobre el punt \(0\) cau al
complement i sobre el punt \(1\), a \(m\). Per això \enquote{\(x\in m^c\)
si i només si \(x\notin m\)} només val sobre punts globals
(observació~\ref{obs:complement-classic}).
\end{description}

\section*{Relacions, imatges, exponencials i parts (capítols 4--7)}

\begin{description}
\item[Un monomorfisme no escindit.] L'única aplicació \(\mathbf0\to\mathbf1\)
és mono i no té cap retracció, perquè no hi ha cap aplicació
\(\mathbf1\to\mathbf0\). L'axioma d'elecció parla dels epimorfismes, no dels
monomorfismes.

\item[Un epimorfisme sense secció.] En els conjunts constants, l'elecció ho
impedeix. Però en els conjunts que varien contínuament sobre un espai (els
feixos sobre la circumferència), la projecció del doble recobriment de la
circumferència sobre ella mateixa és un epimorfisme sense secció: la
\enquote{tria} d'un punt a cada fibra no es pot fer contínuament
\cite[pàg.~14]{bell2006}. Allà l'elecció i la separació global fallen
alhora.

\item[Un element generalitzat d'un exponencial no és una aplicació.] Un
element \(X\to B^A\) és una família d'aplicacions \(A\to B\) indexada per
\(X\); només per a \(X=\mathbf1\) n'és una de sola
(teorema~\ref{teo:punts-exponencial}).

\item[La unió interna sobre elements generalitzats.] Siguin
\(\alpha=\ulcorner1\urcorner\) i \(\beta=\ulcorner0\urcorner\), els noms dels
dos punts com a parts de \(\mathbf2\), i \(a=\id_{\mathbf2}\). Aleshores
\(a\mathrel{\varepsilon}(\alpha\cup\beta)\), perquè la unió és total, però
\(a\) no pertany ni a \(\alpha\) ni a \(\beta\). Les lectures amb \enquote{o}
i \enquote{no} només valen sobre punts
(proposició~\ref{prop:booleanes-puntuals}).
\end{description}

\section*{Nombres naturals, finitud i infinitud (capítols 8--12)}

\begin{description}
\item[El predecessor que no ho és.] La recursora del sistema
\((\N,0,\id_\N)\) compleix \(u\circ s=u\) i \(u\circ0=0\): és la constant
\(0\), no el predecessor. Cal una recursió sobre \(\N\times\N\)
(observació~\ref{obs:predecessor-parell}).

\item[Dues desigualtats que fallen alhora.] Els elements generalitzats
\(m=[0,s\circ0]\) i \(n=[s\circ0,0]\colon\mathbf2\to\N\) no compleixen ni \(m\le n\) ni
\(n\le m\), però el parell pertany a la unió de les dues relacions: la
totalitat de l'ordre és una afirmació sobre parts.

\item[Mono no iso amb la mateixa mida.] El successor \(s\colon\N\to\N\) és
mono i no epi. En els conjunts finits del mateix cardinal, mono, epi i iso
coincideixen; en els infinits, no.

\item[Una unió numerable sense elecció.] Sense l'elecció numerable és
consistent que la recta real sigui una unió numerable de conjunts
numerables (Feferman--Lévy; vegeu~\cite{jech2003}). El
teorema~\ref{teo:unio-numerable} necessita, doncs, l'elecció numerable.

\item[Dues dominacions sense isomorfisme evident.] \(\N\hookrightarrow\N+\mathbf1\)
i \(\N+\mathbf1\hookrightarrow\N\) (pel successor) són monomorfismes oposats,
cap dels dos no és isomorfisme, i tanmateix \(\N\cong\N+\mathbf1\):
Schröder--Bernstein construeix un isomorfisme \emph{nou}
(teorema~\ref{teo:schroder-bernstein}).
\end{description}

\section*{Bons ordres, ordinals i cardinals (capítols 13--19)}

\begin{description}
\item[Rigidesa en un ordre total que no és bo.] A \(\Z\), l'aplicació
\(n\mapsto n-1\) és creixent i injectiva i fa baixar tots els punts; i la
translació és un automorfisme diferent de la identitat. La rigidesa
(lema~\ref{lema:rigidesa}) depèn del mínim de les parts.

\item[Una cadena de bons ordres que no ho és.] Per a cada \(n\),
sigui \(W_n\) la part finita de \(\Z\) formada pels enters entre \(-n\) i
\(0\), amb l'ordre usual. Cada \(W_n\) és un bon ordre i estén l'anterior
\emph{per sota}; formen una cadena la unió de la qual (els enters no positius) no té
mínim. A la prova del lema de Zorn cal que les extensions siguin inicials.

\item[El suprem amb segments tancats.] Amb un sol índex i \(r\) el mínim
d'un bon ordre \(W\), el tipus de \(W_{<r}\) és \(0\) però \(W_{\le r}\) té
tipus \(1\) (observació~\ref{obs:suprem-oberts}).

\item[L'aritmètica ordinal no és commutativa.] \(1+\omega=\omega<\omega+1\)
i \(2\cdot\omega=\omega<\omega\cdot2\); i per l'esquerra la monotonia no és
estricta: \(1<2\) però \(1+\omega=2+\omega\)
(proposició~\ref{prop:monotonia-ordinal}).

\item[La base zero.] Amb la clàusula límit, \(0^\omega\) valdria
\(\sup\{1,0,0,\dots\}=1\), i no \(0\); per això la base \(0\) es tracta a
part (teorema~\ref{teo:exponenciacio-recursio}). En els cardinals, la
monotonia en l'exponent falla per la mateixa raó: \(0^0=1\not\le0=0^1\)
(proposició~\ref{prop:lleis-cardinals}).

\item[Una regla sense unicitat.] Una regla que determina cada \(X_i\) només
llevat d'isomorfisme ---per exemple, \enquote{un conjunt numerable}--- no
permet enganxar les fibres de manera canònica, perquè cada fibra té molts
automorfismes; per això l'esquema SR (principi~\ref{sr:esquema}) demana
isomorfisme únic.
\end{description}
